\documentclass{beamer}

\usepackage[utf8]{inputenc}
\usepackage[french]{babel}
\usepackage{xcolor}
\usepackage{lmodern}
\usepackage{ae, aecompl}

\usepackage{soul}
\usepackage{ulem}
\usepackage{graphicx}
\usepackage{eso-pic}
\usepackage{float}

\usepackage{amssymb}
\usepackage{amsmath}
\usepackage{mathrsfs}
\usepackage{amsthm}
\usepackage{url}
\usepackage{listings}
\usepackage{stmaryrd}
\usepackage{array}
\usepackage{multirow}
\usepackage{verbatim}
\usepackage{hyperref}

\usetheme{Ilmenau}
\setbeamertemplate{footline}
{
  \hbox{\begin{beamercolorbox}[wd=1\paperwidth,ht=2.25ex,dp=1ex,right]{framenumber}%
      \usebeamerfont{framenumber}\insertframenumber{} / \inserttotalframenumber\hspace*{2ex}
    \end{beamercolorbox}}%
  \vskip0pt%
}


\AtBeginSection[]{
  \begin{frame}
    \frametitle{Plan}
    \tableofcontents[currentsection]
  \end{frame}
}

% \AtBeginSubsection{
%   \begin{frame}
%     \frametitle{Plan}
%     \tableofcontents[currentsection,currentsubsection]
%   \end{frame}
% }

\title[Soutenance de Stage]{Soutenance du Stage de deuxième année F4\\Détermination d'un algorithme améliorant l'apprentissage d'un réseau de neurones}
\author{Julien Feuillas}
\institute{\bsc{ISIMA}}
\date{23 Mai 2019}

\begin{document}

\setbeamertemplate{section shaded}[ball unumberded]

\begin{frame}
  \includegraphics[width=0.2\textwidth]{annex/logo_UCA.png}
  \hfill
  \includegraphics[width=0.4\textwidth]{annex/isima.jpeg}
  \titlepage
\end{frame}

% \begin{frame}
%   \begin{itemize}
%     \item bonjour
%     \pause \item c'est une présentation
%   \end{itemize}
%
%   \pause%\includegraphics[width=0.4\textwidth]{annexe/logo_UCA.png}
% \end{frame}

% \begin{frame}{test de blocs}
%   \begin{block}{machin}
%     un truc bleu
%   \end{block}
%
%   \begin{alertblock}{lala}
%     un truc rouge
%   \end{alertblock}
%
%   \begin{exampleblock}{exemple}
%     un exemple vert ?
%   \end{exampleblock}
%
% \end{frame}

\section*{Introduction}

\subsection*{}

\begin{frame}{Objectifs du groupe MMIV}
  \begin{itemize}
    \item MMIV (Mohn Medical Imaging and Visualization Centre)
    \begin{itemize}
      \item Centre de recherche en termes d'imagerie médicale
      \item Basé en Norvège
      \item Dépendant de l'Université de Bergen
      \item Représentante : Mme Renate Grüner
    \end{itemize}
    \item[]
    \item Objectifs
    \begin{itemize}
      \item Mettre en place de nouvelles techniques d'apprentisssage automatique
      \item Création d'offres personnalisées
    \end{itemize}
  \end{itemize}
\end{frame}

\begin{frame}{Travaux réalisés dans le cadre du projet}
  \begin{itemize}
    \item Stage de 6 mois
    \item[]
    \item Création de capteurs
    \begin{itemize}
      \item Utilisation des téléphones portables des clients
      \item Informations recueillies :
      \begin{itemize}
        \item Puissance du signal
        \item Adresse MAC
        \item Date
      \end{itemize}
    \end{itemize}
    \item[]
    \item Mise en place d'une première solution
    \begin{itemize}
      \item Acquisition de données
      \item Algorithme de Machine Learning
    \end{itemize}
  \end{itemize}
\end{frame}

\begin{frame}{Cadre du Projet}
  \begin{itemize}
    \item Nos Objectifs
    \begin{itemize}
      \item Calcul de la position des clients à un instant donné
      \item \'Etude de différentes solutions
    \end{itemize}
    \item[]
    \item Cadre d'étude
    \begin{itemize}
      \item Calcul mathématique direct ?
      \item Machine Learning
      \begin{itemize}
        \item Classification
        \item Régression
      \end{itemize}
    \end{itemize}
  \end{itemize}
\end{frame}

\begin{frame}{Problématique}
  \begin{beamercolorbox}[ht=5ex,dp=1.5ex,center]{frametitle}
    Quelle est la méthode de Machine Learning la plus efficace pour effectuer de la localisation en environnement intérieur ?
  \end{beamercolorbox}
\end{frame}

\section*{Plan}
\begin{frame}{Plan}
  \tableofcontents
\end{frame}

\section[Travaux précédents -- Déroulement]{Travaux précédents et Déroulement du Projet}

\subsection*{}


\begin{frame}{Acquisition des données}
  \begin{itemize}
    \item Découpage de la pharmacie en différentes zones
    \item[]
    \item Méthode d'acquisition des données :
    \begin{itemize}
      \item Téléphone posé dans une zone
      \item \'Emission de signaux pendant une certaine durée
    \end{itemize}
    \item[]
    \item Stockage des données :\\
    \texttt{x;y;date;c1;c2;c3;c4}
  \end{itemize}
\end{frame}

\begin{frame}{Découpage de la pharmacie}
  \begin{center}
    %\includegraphics[scale=0.575]{annexe/rep_pharmacie.pdf}
  \end{center}
\end{frame}

\begin{frame}{Machine Learning}
  \begin{itemize}
    \item Vocabulaire :
    \begin{itemize}
      \item Feature (\textit{c1,c2,c3,c4})
      \item Target (\textit{x,y})
    \end{itemize}
    \item[]
    \item Algorithmes paramétrés
    \item[]
    \item Deux étapes :
    \begin{itemize}
      \item Entraînement
      \item Test
    \end{itemize}
  \end{itemize}
\end{frame}

\begin{frame}{Classification et Régression}
  \begin{itemize}
    \item Classification
    \begin{itemize}
      \item Regroupement en classes
      \item Classes -- Targets
    \end{itemize}
    \item[]
    \item Régression
    \begin{itemize}
      \item Forme de la fonction : features $\mapsto$ targets
      \item Minimisation de l'erreur
      \item Méthodes linéaires :
      \begin{math}
        \left\{
        \begin{array}{c c c c c cccc}
          t_1 &=& \alpha_{1,0} &+& \alpha_{1,1}f_1 &+& \dots &+& \alpha_{1,n}f_n\\
          &&&&\vdots\\
          t_m &=& \alpha_{m,0} &+& \alpha_{m,1}f_1 &+& \dots &+& \alpha_{m,n}f_n
        \end{array}
        \right.
      \end{math}
    \end{itemize}
  \end{itemize}
\end{frame}

\begin{frame}{Outils utilisés}
  \begin{itemize}
    \item Python
    \item[]
    \item Bibliothèques Python :
    \begin{itemize}
      \item Numpy
      \item Scikit-Learn et Pandas
      \item Matplotlib et Seaborn
    \end{itemize}
  \end{itemize}
\end{frame}

\begin{frame}{Déroulement du projet}
  \begin{itemize}
    \item Rendez-vous hebdomadaires
    \item[]
    \item Objectifs définis au cours du projet
  \end{itemize}
\end{frame}

\section[Méthodes]{Présentation des méthodes à implémenter}

\subsection*{}

\begin{frame}{Découpage du jeu de données}
  \begin{center}
    %\includegraphics[scale=0.15]{annexe/kaggle_train_test_split.png}
  \end{center}
\end{frame}

\begin{frame}{Méthodes de Classification}
  \begin{itemize}
    \item K Nearest Neighbors
    \begin{itemize}
      \item Recherche des voisins
      \item Détermination de la classe
    \end{itemize}
    \item[]
    \item Support Vector Machine (SVM)
    \begin{itemize}
      \item Choix d'une méthode
      \item Itération
    \end{itemize}
  \end{itemize}
\end{frame}

\begin{frame}{Métriques pour la Classification}
  \begin{itemize}
    \item $\text{Précision} = \frac{\text{Nombre de prédictions correctes}}{\text{Nombre de données}}$
    \item Matrice de confusion :\\
    \begin{center}
      \begin{tabular}{>{\begin{bf}}c<{\end{bf}} | c c}
        & \textbf{A} & \textbf{B}\\
        \hline
        A & 2345 & 0\\
        B & 213 & 2143
      \end{tabular}
    \end{center}
  \end{itemize}
\end{frame}

\begin{frame}{Régressions linéaire et polynomiale}
  \begin{itemize}
    \item Régression linéaire
    \begin{itemize}
      \item Méthode des moindres carrés
      \item Erreur : $\Vert.\Vert_2^2$
    \end{itemize}
    \item[]
    \item Régression polynomiale
    \begin{itemize}
      \item Paramètre : degré
      \item Exemple :\\
      \begin{math}
        t = a_0+a_1f_1+a_2f_2+a_3f_1^2+a_4f_2^2+a_5f_1f_2
      \end{math}
      \item Méthode linéaire ?
    \end{itemize}
  \end{itemize}
\end{frame}

\begin{frame}{Régressions linéaires paramétrées}
  \begin{itemize}
    \item Régression Ridge
    \begin{itemize}
      \item Redondance d'information entre les individus
      \item Coefficient de pénalité $\alpha$
    \end{itemize}
    \item[]
    \item Régression LASSO
    \begin{itemize}
      \item Redondance d'information dans les features
      \item Coefficient de pénalité $\alpha$
    \end{itemize}
    \item[]
    \item Régression Elastic-Net
    \begin{itemize}
      \item Combinaison des méthodes précédentes
      \item Deux Coefficients :
      \begin{itemize}
        \item $\alpha$ : Coefficient de pénalité
        \item $\rho \in [0,1]$ : Contrôle de la combinaision
      \end{itemize}
    \end{itemize}
  \end{itemize}
\end{frame}

\begin{frame}{Métriques pour la régression}
  \begin{itemize}
    \item Erreur de "distance" :
    \begin{itemize}
      \item Erreur quadratique :
      \begin{itemize}
        \item $MSE(Y,Y_{pred}) = \frac{1}{N_{test}} \sum_{l=1}^{N_{test}} {(Y_l - Y_{pred,l})}^2$
      \end{itemize}
      \item Erreur moyenne absolue
      \begin{itemize}
        \item $MAE(Y,Y_{pred}) = \frac{1}{N_{test}} \sum_{l=1}^{N_{test}} |Y_l - Y_{pred,l}|$
      \end{itemize}
    \end{itemize}
    \item[]
    \item Capacité de prédiction du modèle
    \begin{itemize}
      \item Coefficient de détermination
      \begin{itemize}
        \item $R^2(Y,Y_{pred}) = 1 - \frac{\sum_{l=1}^{N_{test}} (Y_l - Y_{pred,l})^2}{\sum_{l=1}^{N_{test}} (Y_l - \bar{Y})^2}$
      \end{itemize}
      \item Score de variance expliquée
      \begin{itemize}
        \item $EVS(Y,Y_{pred}) = 1 - \frac{Var(Y - Y_{pred,l})}{Var(Y)}$
      \end{itemize}
    \end{itemize}
  \end{itemize}
\end{frame}

\begin{frame}{Efficacité des méthodes de régression}
  \begin{itemize}
    \item Résultats dépendants du découpage initial
    \item[]
    \item \'Ecriture d'une fonction Python :
    \begin{itemize}
      \item \texttt{apply\_regressions}
      \item Plusieurs applications du même modèle de régression
      \item Renvoie la valeur des métriques pour chaque application
      \item Variables d'entrée :
      \begin{itemize}
        \item Nombre de régressions
        \item Le modèle à appliquer
        \item Conservation des lignes dupliquées ?
      \end{itemize}
    \end{itemize}
    \item[]
    \item Utilisable pour la détermination de paramètres
  \end{itemize}
\end{frame}

\section{Résultats}

\subsection*{}

\begin{frame}{K Nearest Neighbors}
  \begin{center}
    %\includegraphics[scale=0.57]{annexe/tables/K_Nearest/K.pdf}
  \end{center}
  \begin{itemize}
    \item Précision = 98\%
  \end{itemize}
\end{frame}

\begin{frame}{Support Vector Classification}
  \begin{center}
    %\includegraphics[scale=0.57]{annexe/tables/SVC/svc.pdf}
  \end{center}
  \begin{itemize}
    \item Précision = 68\%
  \end{itemize}
\end{frame}

\begin{frame}{Détermination des paramètres : polynomiale}
  \begin{center}
    %\includegraphics[scale=0.57]{annexe/param_poly.pdf}
  \end{center}
\end{frame}

\begin{frame}{Détermination des paramètres : Ridge}
  \begin{center}
    %\includegraphics[scale=0.57]{annexe/param_ridge.pdf}
  \end{center}
\end{frame}

\begin{frame}{Détermination des paramètres : LASSO}
  \begin{center}
    %\includegraphics[scale=0.57]{annexe/param_lasso.pdf}
  \end{center}
\end{frame}

\begin{frame}{Détermination des paramètres : Elastic-Net}
  \begin{itemize}
    \item Méthode précédente non utilisable
    \item[]
    \item Utilisation d'un objet déjà implémenté :
    \begin{itemize}
      \item \texttt{MultiTaskElasticNetCV} de \texttt{sklearn.linear\_model}
      \item Validation croisée + méthode de régression
    \end{itemize}
    \item[]
    \item Paramètre de construction :
    \begin{itemize}
      \item Nombre de plis à effectuer
      \item Un tableau de valeurs possibles pour $\alpha$ er $\rho$
    \end{itemize}
  \end{itemize}
\end{frame}

\begin{frame}{Tableau de résultats}
  \begin{center}
    \begin{tabular}{|c|c|c|c|c|c|}
      \hline
       & Duplications ? & MSE & MAE & R2 & EVS\\
      \hline
      \multirow{2}*{Linéaire} & oui & 1.862 & 1.097 & 0.642 & 0.643\\
      \cline{2-6}
      & non & 2.061 & 1.158 & 0.602 & 0.606\\
      \hline
      \multirow{2}*{Polynomiale} & oui & 0.900 & 0.666 & 0.824 & 0.825\\
      \cline{2-6}
      & non & 1.299 & 0.773 & 0.717 & 0.720\\
      \hline
      \multirow{2}*{Elastic-Net} & oui & 1.864 & 1.098 & 0.642 & 0.642 \\
      \cline{2-6}
      & non & 2.064 & 1.160 & 0.601 & 0.604\\
      \hline
    \end{tabular}
  \end{center}
\end{frame}

\begin{frame}{Comparaison entre classification et régression}
  \begin{itemize}
    \item Efficacité de la régression :
    \begin{itemize}
      \item Au sens des métriques de régression
      \item Différentes des métriques de classification
    \end{itemize}
    \item[]
    \item Passage régression $\to$ classification
    \begin{itemize}
      \item Pour la régression polynomiale
    \end{itemize}
  \end{itemize}
\end{frame}

\begin{frame}{Passage de la régression à la classification}
  \begin{center}
    %\includegraphics[scale=0.57]{annexe/assoc_reg_class.pdf}
  \end{center}
\end{frame}

\begin{frame}{Résultats de la régression polynomiale pour la classification}
  \begin{itemize}
    \item Résultats :
    \begin{itemize}
      \item Précision : 54\%
      \item Termes extra-diagonaux de la matrice de confusion
    \end{itemize}
    \item[]
    \item Résultats peu précis
    \begin{itemize}
      \item Zones trop grandes
      \item Nécessité d'un nouveau jeu de données
    \end{itemize}
  \end{itemize}
\end{frame}

\begin{frame}{Graphique de densité}
  \begin{center}
    %\includegraphics[scale=0.57]{annexe/comp_densite.pdf}
  \end{center}
\end{frame}

\section*{Conclusion}

\begin{frame}{Résumé du travail réalisé}
  \begin{itemize}
    \item Méthodes de régression implémentées
    \begin{itemize}
      \item Méthodes linéaires
      \item Méthodes probabilistes ? Régressions à noyau ?
    \end{itemize}
    \item[]
    \item Résultats obtenus
    \begin{itemize}
      \item Non concluants
      \item Méthodes employées réutilisables
    \end{itemize}
    \item[]
    \item Ce qu'il reste à faire
    \begin{itemize}
      \item Mise en place de nouvelles régressions
      \item \'Etude d'un nouveau jeu de données
    \end{itemize}
  \end{itemize}
\end{frame}

\section*{}

\begin{frame}{Remerciements}
  \begin{beamercolorbox}[ht=2.5ex,dp=1.5ex,center]{frametitle}
    Nous vous remercions pour votre attention
  \end{beamercolorbox}
\end{frame}

\appendix

\section*{Annexe}

% \begin{frame}{Triangulation}
%   \begin{center}
%     %\includegraphics[scale=0.57]{annexe/satellites.jpg}
%   \end{center}
% \end{frame}
%
% \begin{frame}{Exemple de Classification}
%   \begin{center}
%     %\includegraphics[scale=0.5]{annexe/ex_classif.pdf}
%   \end{center}
% \end{frame}
%
% \begin{frame}{Exemple de Régression}
%   \begin{center}
%     %\includegraphics[scale=0.5]{annexe/ex_regression.pdf}
%   \end{center}
% \end{frame}
%
% \begin{frame}{Plan de la pharmacie}
%   \begin{center}
%     %\includegraphics[scale=0.5]{annexe/plan_pharmacie.png}
%   \end{center}
% \subsection*{}
%
% \end{frame}
%
% \begin{frame}{Underfitting et Overfitting}
%   \begin{center}
%     %\includegraphics[scale=0.57]{annexe/underoverfitting.pdf}
%   \end{center}
% \end{frame}

\end{document}
